% Formen-Rechner: Dokumentation. TODO: Beispiele ergaenzen
% FIXME XXX NOTE
\documentclass[11pt,a4paper]{article}
\usepackage[utf8]{inputenc}
\usepackage{amsmath, amssymb, graphicx, listings, hyperref, verbatim}
\input{kopf}
\include{anhang}
\newcommand{\flaeche}[2][r]{\pi #1^{2} + #2}
\renewcommand{\arraystretch}{1.2}
\newenvironment{hinweis}[1]{\begin{quote}\textbf{#1}: }{\end{quote}}
\DeclareMathOperator{\umfang}{U}
\def\firma{Muster \& Co}
\setlength{\parindent}{0pt}

\title{Der Formen-Rechner}
\author{Team \and G\"ast\'e \c{c} \ss{} \ae \O}
\date{\today}

\begin{document}
\maketitle
\tableofcontents

\begin{abstract}
  Berechnet Fl\"ache und Umfang. Preis: 10\,\% oder \$5 -- siehe unten --- fertig~sofort.
\end{abstract}

\part{Grundlagen}
\chapter*{Vorwort}
\section{Einleitung}\label{sec:einleitung}
\subsection*{Ziele}
\subsubsection{Details}
\paragraph{Absatz} Text mit \emph{Betonung}, \textit{kursiv}, \textbf{fett \textit{beides}}
und \texttt{Schreibmaschine}. {\bfseries Fett} {\em alt} {\Large gross}\\[2ex]
Siehe Abschnitt~\ref{sec:einleitung} auf Seite~\pageref{sec:einleitung}, \eqref{eq:kreis}
und \cite[S.~5]{knuth84,lamport94}. \nocite{*} \footnote{Eine Fussnote.}
\url{https://example.org/a_b#c} \href{https://example.org}{Webseite}

\section{Mathematik}
Inline: $A = \pi r^2$ und \(U = 2\pi r\), $x_{i+1} \leq \alpha \cdot \beta - \sqrt{\gamma}$.
Abgesetzt: $$\sum_{i=0}^{n} i = \frac{n(n+1)}{2}$$
\[
  \int_0^\infty e^{-x^2}\,\mathrm{d}x = \frac{\sqrt{\pi}}{2} \quad \text{fuer alle $x$}
\]
\begin{equation}\label{eq:kreis}
  A(r) = \pi r^{2}, \qquad \left( \frac{a}{b} \right) \Rightarrow \left\{ x \in \mathbb{R} \right.
\end{equation}
\begin{align*}
  a &= b + c & d &\neq e \\
  f &\approx \lim_{x \to 0} \sin(x)/x % Kommentar in Mathe
\end{align*}
\begin{eqnarray}
  y & = & \begin{pmatrix} 1 & 0 \\ 0 & 1 \end{pmatrix} \nonumber
\end{eqnarray}
\begin{displaymath} \vec{v} = (v_1, \ldots, v_n)^T \end{displaymath}
\begin{math} \Omega \end{math}

\section{Listen und Tabellen}
\begin{itemize}
  \item Kreis
  \item[--] Rechteck
\end{itemize}
\begin{enumerate}
  \item erstens \item zweitens
\end{enumerate}
\begin{description}
  \item[Form] eine Fl\"ache
\end{description}
\begin{table}[htbp]
  \centering
  \begin{tabular}{|l|c|r@{.}l|}
    \hline
    Form & Seiten & \multicolumn{2}{c|}{Wert} \\ \cline{1-2}
    Kreis & 0 & 3 & 14 \\
    \hline
  \end{tabular}
  \caption{Formen}\label{tab:formen}
\end{table}
\begin{figure}
  \includegraphics[width=0.5\textwidth]{bild.png}
  \caption[kurz]{Ein Bild mit $r = 1$.}
\end{figure}

\section{Woertlich}
Mit \verb|\foo{$x$}| und \verb*+a b+ im Text.
\begin{verbatim}
  def flaeche(r):   % kein Kommentar $x$
      return \pi
\end{verbatim}
\begin{lstlisting}[language=Python]
print("Hallo %s" % name)  # kein TeX \section{x}
\end{lstlisting}
\begin{comment}
  Alles Kommentar. TODO
\end{comment}
\iffalse versteckt \fi
%begin-include
Eingebunden.
%end-include

\makeatletter
\def\@interner@befehl#1#2{#1\relax}
\makeatother
\ifx\foo\undefined \else \fi
\begin{hinweis}{Achtung} Klammern { verschachtelt { tief } } fertig. \end{hinweis}
\bibliographystyle{plain}
\bibliography{literatur}
\appendix
\end{document}
Nach dem Ende. \section{ignoriert}
